\section{Conclusion}
\label{ch:conclusion}

This thesis examines how effectively the chatbot answers real institutional questions and provides useful guidance. Its reference architecture describes the path from heterogeneous sources through retrieval and generation to delivery, while the empirical study tests a bounded part of that path. The contribution is an inspectable account of the system together with a controlled evaluation of assistance and retrieval policy.

The main evaluation covers 96 historical email questions and 96 questions from selectively recorded production logs, following a separate 20-case pilot. Under a shared retrieval threshold without fallback, the retained AI labels identify supported answering or useful routing in 31/96 email outcomes and 53/96 production outcomes, with useful routing contributing substantially to both. Incomplete responses and unsupported material claims remain visible weaknesses.

Lowering the document threshold changes the combined totals to 30/96 and 52/96, a paired difference of $-1.04$ percentage points in each sample with intervals spanning zero. The intervention substantially increases document context, and improves ranked retrieval in the component study, without a detectable assistance benefit. The fallback comparisons also show no clear benefit, although limited exposure constrains their interpretation. A second, forty-question comparison on an expanded document collection tested a prompt that checks source applicability and preserves procedural steps; with the full generation evidence available to the assessor, it changed strict completion by $-2.5$ points.

These results rest on AI assessments whose limits the audit makes explicit. Empty fact sets can support conversational labels without establishing factual accuracy, references, generation and primary judgments share a model family, and repeated and cross-model assessments disagree. The confidence intervals describe sampling variation conditional on the retained labels and cannot absorb these systematic uncertainties.

The reference architecture made these results interpretable. By separating source preparation, evidence selection, answer construction and delivery, it allowed the observed failures to be attributed to stages: the threshold change largely repaired evidence selection, while incomplete answers and unsupported claims, which belong to answer construction, persisted. For the studied system, completeness checks, explicit clarification and source authority metadata are therefore more promising than admitting more retrieved text. The next step is task-specific reference review with explicit authority checks, human assessment of label reliability and a newly frozen held-out sample.

The supported conclusion is that the chatbot can receive positive assistance labels on selected institutional questions, yet its effectiveness remains limited by incomplete answers, unsupported claims and uncertainty in assessment. Neither more retrieved text nor the tested prompt change establishes better assistance. A stronger evaluation must make the expected task, supporting evidence and appropriate response behavior explicit, so that a complete factual answer, a useful clarification and a necessary referral can each be assessed on their own terms.
